\documentclass{beamer}
\usetheme{Warsaw}
\usebackgroundtemplate{\includegraphics[width=\paperwidth]{format/libresoft-bg-soft.png}}
\usepackage[spanish]{babel}
\usepackage[latin1]{inputenc}
\usepackage{graphics}
\usepackage{amssymb} % Simbolos matematicos

%\definecolor{libresoftgreen}{RGB}{162,190,43}
%\definecolor{libresoftblue}{RGB}{0,98,143}

%\setbeamercolor{titlelike}{bg=libresoftgreen}

%% Metadatos del PDF.
\hypersetup{
  pdftitle={Evaluation of Libre Software Projects / Master on Libre Software (URJC)},
  pdfauthor={Daniel Izquierdo Cortazar},
  pdfcreator={GSyC/LibreSoft, Universidad Rey Juan Carlos},
  pdfproducer=PDFLaTeX,
  pdfsubject={Evaluation},
}
%%

%\includeonly{presentation}
%\includeonly{introduction}
%\includeonly{history}
%\includeonly{consequences}
%\includeonly{public_administrations}
%\includeonly{trials,myths}

\AtBeginSection[]
{
\begin{frame}<beamer>
\begin{center}
{\Huge \insertsection}
\end{center}
\end{frame}
}

\begin{document}

\title{Evaluation of Libre Software Projects}
\subtitle{Master on Libre Software (URJC) \\
\url{http://master.libresoft.es}}
\author{Daniel Izquierdo-Cortazar}
\institute{dizquierdo@libresoft.es \\
GSyC/LibreSoft, Universidad Rey Juan Carlos}

\date{November 2011}

\frame{
\maketitle
\begin{center}
\includegraphics[width=6cm]{format/gsyc-urjc}
\end{center}
}


% Si el titulo o el autor se quieren acortar para los pies de p�gina
% se pueden redefinir aqu�:
%\title{Titulo corto}
%\author{Autores abreviado}


%% LICENCIA DE REDISTRIBUCION DE LAS TRANSPAS
\frame{
~
\vspace{3cm}

\begin{flushright}
\copyright 2011 Daniel Izquierdo Cortazar. \\

Some rights reserved. \\
This document is distributed under the \\
Creative Commons Attribution-ShareAlike 3.0 licence, \\
available in \\
\url{http://creativecommons.org/licenses/by-sa/3.0}

The original version of this document is available at \\
\url{http://master.libresoft.es}
\end{flushright}
}
%%



%% presentation.tex
%%
%% Presentation of the course ``Introduction'' of the Official Master on Libre Software (URJC)
%% http://master.libresoft.es
%%

%%---------------------------------------------------------------------
%%---------------------------------------------------------------------

\section{Introduction to Evaluation of FLOSS Projects}

%%---------------------------------------------------------------

\begin{frame}
\frametitle{Why is this important?}

\begin{itemize}
\item The evaluation of projects is a critical task for IT managers
\item The costs of selecting the wrong technology could be later unapproachable
\end{itemize}
\end{frame}



%%---------------------------------------------------------------

\begin{frame}
\frametitle{General questions made}

\begin{itemize}
\item Which software fits best my current requirements?
\item Requirements such as:
 \begin{itemize}
  \item Functional requirements: looking for specific functionality
  \item Technical requirements: looking for specific technologies
 \end{itemize}
\end{itemize}
\end{frame}

%%---------------------------------------------------------------

\begin{frame}
\frametitle{More requirements}

\begin{itemize}
\item Safety, security, reliability, robustness, testability
\item adaptability, complexity, maintainability, portability

\end{itemize}
\end{frame}

%%---------------------------------------------------------------

\begin{frame}
\frametitle{Even more}

\begin{itemize}
\item Reusability, installability, efficiency,
\item correctness, modularity
\end{itemize}
\end{frame}

%%---------------------------------------------------------------

\begin{frame}
\frametitle{And FLOSS projects?}

\begin{itemize}
\item Is the license included above?
\item Social networks: territoriality, hubs of people (or knowledge)
\item Others...
\end{itemize}
\end{frame}

%%---------------------------------------------------------------

\begin{frame}
\frametitle{Selecting proprietary systems}

\begin{itemize}
\item Some limitations due to depend on proprietary software:
 \begin{itemize}
  \item Cost: generally higher
  \item Closed source code: some requirements could not been measured: e.g. quality of the software
  \item There is not influence in the roadmap at all
  \item Formats
  \item Others
 \end{itemize}
\end{itemize}
\end{frame}

%%---------------------------------------------------------------

\begin{frame}
\frametitle{Selecting FLOSS systems}

\begin{itemize}
\item Some limitations when selecting FLOSS projects:
 \begin{itemize}
  \item Limitless options
  \item In some cases there is not professional support (as requested by companies)
  \item Longevity (although this is also a problem in proprietary software)
  \item Stability: following a ``release often, release early'' may be seen as a weak point by companies
 \end{itemize}
\end{itemize}
\end{frame}

%%---------------------------------------------------------------

\begin{frame}
\frametitle{Thus...}

\begin{itemize}
\item If you need a specific technology then it is clear the option
\item However it is not clear if there are some potential products that fit your requirements
\item How to select them?
\end{itemize}
\end{frame}

%%---------------------------------------------------------------
\section{Initial Assessment}


%%---------------------------------------------------------------

\begin{frame}
\frametitle{FLOSS Assessments}

\begin{itemize}
\item Light-weight: OpenBRR or QSoS
\item Heavy-weight: QualOSS
\item Those will be studied in following lessons
\end{itemize}
\end{frame}


%%---------------------------------------------------------------

\begin{frame}
\frametitle{FLOSS Assessments}

\begin{itemize}
\item However, what type of data do they retrieve?
 \begin{itemize}
   \item Programming language
   \item Quality of the source code
   \item General architecture
   \item License
   \item Specific functionality
   \item Support
 \end{itemize}
\end{itemize}
\end{frame}


%%---------------------------------------------------------------

\begin{frame}
\frametitle{How to proceed?}

\begin{itemize}
\item Following specific methodologies that facilitate this process (e.g.: Goal-Question-Metric).
\item Following specific quality models (that define methodologies)
\end{itemize}
\end{frame}

\section{Exercises}

\begin{frame}
\frametitle{Exercise}

\begin{itemize}
\item Goal: We have to install 5,000 FLOSS e-mail clients in a company. Specific development and
translations will be needed in any case. We will re-sell the product and provide support.
\item Specific mission: Select five requirements that you think are important when selecting this e-mail client.
\item Other considerations: This should be made based on your current and personal experience.
\item Methodology: Create an initial methodology to be followed for other people.
\end{itemize}
\end{frame}


\begin{frame}
\frametitle{Exercise}

\begin{itemize}
\item Improve the previous exercise with a GQM approach
\end{itemize}
\end{frame}


\begin{frame}
\frametitle{Exercise}

\begin{itemize}
\item Improve the previous exercise adding info from the ISO 9126 that could be useful
\end{itemize}
\end{frame}


\end{document}